\section{5. Conclusion}


In this paper, we propose a novel Unified Dual-Mask Physical Model for non-Lambertian light field depth estimation, addressing the fundamental challenge where conventional photo-consistency and epipolar geometry assumptions fail on complex reflectances. To bridge the gap between data-driven representation and physical optics, our approach systematically integrates material-specific physical priors into the depth inference pipeline. By decoupling the light field into distinct reflectance components and applying physics-driven feature modulation, the proposed framework effectively circumvents the depth estimation degradation caused by specularities, glossiness, and translucency, thereby achieving robust and highly accurate geometry recovery in complex real-world scenes.

The core contributions of this work are summarized as follows:
\textit{ \textbf{MRI-Inspired Angular Frequency Domain Material Parsing:} We introduce a novel material parsing mechanism based on angular frequency domain analysis, drawing inspiration from MRI techniques. This approach overcomes the limitations of traditional spatial texture reliance and black-box CNN classifiers (which yielded merely 24.6\% accuracy in baseline evaluations), enabling efficient and highly interpretable material property extraction from a physical optics perspective. It provides exceptionally reliable physical priors for depth estimation in non-Lambertian scenarios.
} \textbf{Dual-Mask Modulated Physics-Driven Feature Representation:} We develop a physics-driven light field feature expression model utilizing dual-mask modulation. By transforming conventional implicit, data-driven feature extraction into explicit, physically meaningful mask modulation, this mechanism significantly enhances the representation capability for mixed materials and complex illumination conditions. It effectively alleviates depth estimation failures caused by texture copying in non-Lambertian regions, substantially improving model interpretability.
\textbf{Component-Aware Multi-Strategy Adaptive Depth Estimation and Fusion:} We design a component-aware adaptive framework that fundamentally resolves the complete failure of traditional Epipolar Plane Image (EPI) methods in non-Lambertian regions due to violated view-consistency assumptions. This multi-strategy fusion mechanism enables a single unified model to process Lambertian, non-Lambertian, and mixed urban scenes with high precision, drastically boosting robustness and generalization capabilities in complex real-world environments.

Ultimately, this work marks a significant step forward in transitioning light field depth estimation from purely empirical, data-driven paradigms to physically grounded, interpretable models. By successfully unifying the geometric inference of diverse material properties under a single physical framework, our method not only establishes a new state-of-the-art for non-Lambertian light field analysis but also lays a solid theoretical and practical foundation for advanced downstream 3D vision tasks. We anticipate that the proposed physical priors and unified modeling strategies will broadly benefit and accelerate applications in autonomous navigation, robotic manipulation, and immersive 3D reconstruction within unconstrained, real-world environments.


The transition from spatial-domain convolution to angular frequency-domain analysis resolves a fundamental ill-posedness inherent in limited-baseline light field sampling. In a $9 \times 9$ angular patch, spatial CNNs conflate high-frequency spatial textures with specular highlights due to the lack of sufficient receptive field and angular baseline, resulting in the observed 24.6\% baseline accuracy. By mapping the angular sampling to the $k$-space via 2D-DFT, we effectively decouple the material reflectance from spatial albedo. From a signal processing standpoint, Lambertian surfaces manifest as a concentrated energy peak at the DC component in the angular frequency domain, whereas specular and scattering components exhibit distinct spectral broadening governed by the microfacet distribution function. This spectral separation ensures that the material parsing is invariant to spatial texture variations, providing a mathematically well-posed foundation for extracting physical priors.

Furthermore, the integration of the medium mask and angular direction mask fundamentally alters the optimization landscape of the Radiance Tensor Field (RTF). In conventional formulations, non-Lambertian reflections inflate the rank of the RTF, rendering low-rank matrix completion or factorization highly ambiguous. The proposed dual-mask modulation acts as an explicit rank regularizer. By modulating the RTF $\mathcal{T}$ with the physically derived roughness parameter $r = a/\lambda$ and the deflection vector field, the high-rank non-Lambertian tensor is projected into a lower-dimensional manifold. This process can be formulated as a constrained rank-minimization problem:
\begin{equation}
\label{eq:eq34}
\min_{\mathbf{D}} \|\mathcal{T} \odot (\mathbf{M}_{medium} \otimes \mathbf{M}_{angular}) - \mathbf{D}\|_F^2 \quad \text{s.t.} \quad \text{rank}(\mathbf{D}) \le k_{phys}
\end{equation}
where $\mathbf{M}_{medium}$ and $\mathbf{M}_{angular}$ denote the dual masks, and $k_{phys}$ is the physically constrained rank bound. This explicit physical factorization significantly reduces the solution space ambiguity during depth regression, ensuring that the geometric disparities are not corrupted by the high-frequency angular variations of glossy surfaces, thereby stabilizing the epipolar plane image (EPI) (EPI) slope estimation even in severely textureless regions. Consequently, the depth inference pipeline transitions from a purely data-driven heuristic search to a physics-constrained optimization, drastically reducing the reliance on massive annotated datasets for non-Lambertian edge cases.

Despite the robust performance, the proposed physical model operates under specific boundary conditions that warrant further investigation. First, the efficacy of the MRI-inspired 2D-DFT relies on the Nyquist-Shannon sampling theorem in the angular domain. When the angular resolution degrades to $5 \times 5$ or $3 \times 3$, severe spectral aliasing occurs in the $k$-space, causing the energy leakage of specular lobes into the mid-frequency scattering bands. Future work could explore super-resolution in the angular frequency domain to mitigate this aliasing. Second, the current medium mask formulation assumes single-bounce surface interactions. For highly translucent materials exhibiting complex subsurface scattering, the photon transport involves multiple bounces, which violates the explicit directional deflection assumption and necessitates a volumetric rendering extension or a diffusion-based physical approximation to accurately model the light transport. Moreover, dynamic scenes involving non-rigid deformations may introduce temporal inconsistencies in the angular frequency signatures, requiring the integration of temporal smoothing filters in the $k$-space.

From an engineering deployment perspective, while the 2D-DFT material parsing maintains an $O(HW)$ computational complexity, the dual-mask modulation introduces non-trivial memory access patterns. The pixel-wise tensor modulation requires frequent off-chip memory fetches for the RTF and mask tensors, potentially creating a memory bandwidth bottleneck on edge devices. Optimizing the memory hierarchy through operator fusion and leveraging specialized hardware accelerators for complex-valued frequency operations will be critical for real-time autonomous navigation applications. Additionally, quantizing the frequency-domain features without losing the subtle phase information of the scattering lobes remains an open challenge for integer-arithmetic neural processing units (NPUs), highlighting a crucial direction for future hardware-software co-design.

Despite the demonstrated efficacy of the proposed unified dual-mask physical model, several limitations remain that warrant further investigation from a signal processing and optical engineering perspective. First, the reliability of the angular frequency domain analysis is inherently constrained by the angular sampling density. The current formulation relies on a dense $9 \times 9$ angular grid to execute the two-dimensional discrete Fourier transform (2D-DFT). In scenarios involving sparse angular sampling, such as $3 \times 3$ or $5 \times 5$ micro-lens arrays, the frequency-domain resolution degrades significantly. This sparsity introduces spectral aliasing, causing the high-frequency energy spread of specular reflections to overlap with the mid-frequency bands of scattering materials, thereby compromising the fidelity of the extracted physical priors.

Second, the dual-mask modulation mechanism assumes predominantly single-bounce light-matter interactions. While the medium and angular direction masks effectively decouple local reflectance properties, they do not explicitly account for complex global illumination effects, such as inter-reflections in concave geometries or deep subsurface scattering in translucent media. In such regions, the radiance tensor field (RTF) rank analysis may deviate from theoretical low-rank assumptions, as multi-bounce radiosity introduces non-linear angular correlations that the current linear mask modulation cannot fully disentangle.

To address these limitations, future work will focus on three primary directions. First, we plan to develop an adaptive frequency-domain zero-padding and interpolation strategy guided by analytical BRDF priors. This will mitigate spectral aliasing in sparse light fields, extending the robustness of the material parsing mechanism to consumer-grade plenoptic cameras. Second, we aim to incorporate global illumination models into the physical feature representation. By integrating multi-bounce rendering equations, the dual-mask framework can be generalized to handle severe inter-reflections and complex volumetric scattering. Finally, we will explore hardware-software co-design paradigms. Implementing the angular 2D-DFT and mask modulation directly in the optical domain using programmable metasurfaces or neuromorphic event-based sensors could drastically reduce computational overhead, paving the way for real-time, low-power non-Lambertian depth estimation in autonomous systems.